% 哥白尼原则（综述）
% license CCBYSA3
% type Wiki

本文根据 CC-BY-SA 协议转载翻译自维基百科\href{https://en.wikipedia.org/wiki/Copernican_principle}{相关文章}。

在物理宇宙学中哥白尼原理（Copernican principle）指出：人类并非宇宙中的特权观察者，\(^\text{[1]}\)来自地球的观测结果可视为来自宇宙中平均位置的代表性观测。该原理以哥白尼的日心说（heliocentrism）命名，是一种工作假设，其思想源自对哥白尼“地球运动”论证的宇宙学延伸。\(^\text{[2]}\)
\subsection{起源与意义}
\begin{figure}[ht]
\centering
\includegraphics[width=6cm]{./figures/3229b0a8fd43c8a4.png}
\caption{图 “M”（取自拉丁语Mundus，意为“世界”），出自约翰内斯·开普勒（Johannes Kepler）于 1617–1621 年间撰写的《哥白尼天文学概要》（Epitome Astronomiae Copernicanae），展示了地球只是众多类似恒星体系中的一个。} \label{fig_CoPr_1}
\end{figure}
赫尔曼·邦迪（Hermann Bondi）在二十世纪中期将该原理命名为“哥白尼原理”，但该思想本身可追溯至十六至十七世纪期间从托勒密体系（Ptolemaic system）向新范式的转变——后者不再将地球置于宇宙中心。哥白尼提出，行星的运动可通过假设太阳处于静止的中心位置来解释，从而与地心说形成对比。他认为，行星的视逆行运动是一种由地球绕太阳运动所造成的幻象，而在哥白尼模型中，太阳被置于宇宙的中心。值得注意的是，哥白尼的主要动机是出于对旧体系的技术性不满，而非出于支持任何“平庸原理”（principle of mediocrity）的哲学立场。\(^\text{[3]}\)

尽管哥白尼的日心模型常被描述为“将地球降级”——使其失去在托勒密地心模型中的中心地位，但真正采纳这一新宇宙观的，是哥白尼的继承者，尤其是十六世纪的焦尔达诺·布鲁诺（Giordano Bruno）。在传统观念中，地球的中心位置被认为是“最低且最污秽之处”；而伽利略则指出，地球实际上参与了“群星的舞蹈”，而非停留在“宇宙污秽与浮渣的积聚之所”。\(^\text{[4][5]}\)

至二十世纪后期，卡尔·萨根（Carl Sagan）以诗意的语言进一步强化了哥白尼原理的精神，他问道：“我们是谁？我们发现自己居住在一颗微不足道的行星上，这颗行星环绕着一颗平凡的恒星，处于银河系的某个被遗忘的角落，而宇宙中星系的数量远远超过人类的数量。”\(^\text{[6]}\)

尽管哥白尼原理（Copernican principle）起源于对过去假设的否定——如地心说（geocentrism）、日心说（heliocentrism）或银河中心说（galactocentrism）等，这些假设都认为人类处于宇宙的中心——但哥白尼原理的内涵比非中心论（acentricism）更为深刻。非中心论仅仅主张人类不位于宇宙中心，而哥白尼原理在承认非中心论的前提下，更进一步断言：人类观察者，或来自地球的观测结果，是宇宙中“平均位置”观测的代表。迈克尔·罗文-罗宾逊（Michael Rowan-Robinson）强调，哥白尼原理是现代思想的门槛性检验标准。他指出：“显而易见，在人类历史的后哥白尼时代，没有一个见多识广且理性的人会认为地球在宇宙中占据着独一无二的位置。”\(^\text{[7]}\)

现代宇宙学的大部分理论建立在一个前提之上：宇宙学原理（cosmological principle）在最大尺度上几乎成立，但并非绝对成立。哥白尼原理代表了支持这一前提所需的、不可约的哲学假设——当其与观测结果结合时，构成了宇宙学的逻辑基础。如果假定哥白尼原理成立，并且从地球的视角观测到宇宙在各个方向上都是各向同性的（isotropic，即在各个方向上看起来相同），则可以推断宇宙在总体上是均匀的（homogeneous，即在任意给定时刻处处相似），并且围绕任意一点都是各向同性的。这两种条件共同构成了宇宙学原理。\(^\text{[7]}\)

实际上，天文学家观测到的宇宙在星系超星系团（galactic superclusters）、纤维状结构（filaments）以及巨大空洞（great voids）等尺度上表现出明显的不均匀性。在当代主流宇宙学模型——即Lambda-CDM 模型（$\Lambda$CDM model）中，宇宙被预测为在更大的观测尺度上将逐渐趋于均匀和各向同性，当观测尺度超过约 $260$ 百万秒差距（$\mathrm{Mpc}$）时，可探测到的结构将极少。\(^\text{[8]}\)然而，来自星系团、\(^\text{[9][10]}\)类星体（quasars）、\(^\text{[11]}\)以及 Ia 型超新星（type Ia supernovae）\(^\text{[12]}\)的最新证据表明，宇宙在大尺度上可能违反了各向同性。此外，科学家还发现了多种大尺度结构，如 Clowes–Campusano LQG、斯隆长城（Sloan Great Wall）、U1.11、巨大类星体群（Huge-LQG）、武仙–北冕座大墙（Hercules–Corona Borealis Great Wall）、巨弧（Giant Arc）以及本地空洞（Local Hole），\(^\text{[13][14][15][16]}\)这些结构都暗示着宇宙的均匀性可能被破坏。

在与可观测宇宙半径相当的尺度上，我们观察到从地球出发的距离越大，宇宙结构呈现出系统性的变化。例如，在更遥远的距离处，星系中包含更多的年轻恒星，星系间的聚集程度更低，而类星体（quasars）的数量则显得更多。如果假定哥白尼原理（Copernican principle）成立，那么这意味着这些观测结果是宇宙随时间演化的证据：这些遥远的光线花费了几乎整个宇宙年龄才到达地球，因此展示的是宇宙在早期的样貌。所有可观测到的最遥远的光——宇宙微波背景辐射（cosmic microwave background radiation）——在至少千分之一的精度范围内是各向同性的。

邦迪（Bondi）和托马斯·戈尔德（Thomas Gold）利用哥白尼原理提出了完美宇宙学原理（perfect cosmological principle），该原理主张宇宙在时间上也是均匀的（homogeneous in time），并由此成为稳态宇宙学（steady-state cosmology）的理论基础。c然而，这一观点与前述关于宇宙学演化的观测证据严重冲突：宇宙自大爆炸（Big Bang）以来经历了极其不同的物理条件，并将在暗能量（dark energy）影响不断增强的作用下继续向极端不同的状态演化——显然趋向于“热寂”（Big Freeze）或“撕裂”（Big Rip）结局。

自二十世纪九十年代以来，该术语（“哥白尼原理”）也被用于（与“哥白尼方法”Copernicus method可互换）J. Richard Gott 提出的基于贝叶斯推断（Bayesian inference）的事件持续时间预测模型中，该模型是“世界末日论证”（Doomsday argument）的推广形式。\,\(^\text{[clarification needed]}\)
\subsubsection{历史背景}
在“哥白尼原理”（Copernican principle）这一术语被正式提出之前，许多以往的假设——例如地心说（geocentrism）、日心说（heliocentrism）以及银河中心说（galactocentrism）——都认为地球、太阳系或银河系分别位于宇宙的中心。然而，这些假设均被证明是错误的。哥白尼革命（Copernican Revolution）使地球“退位”，成为围绕太阳运转的众多行星之一。哈雷（Halley）首次提及了恒星自行（proper motion）的现象。威廉·赫歇尔（William Herschel）发现太阳系在我们盘状的银河系中穿行运动。埃德温·哈勃（Edwin Hubble）进一步表明，银河系只是宇宙中无数星系之一。对银河系在宇宙中的位置与运动的研究，最终导致了大爆炸理论（Big Bang theory）的提出，并奠定了整个现代宇宙学的基础。
\subsubsection{现代检验}
近期以及计划中的与宇宙学原理（cosmological principle）和哥白尼原理相关的实验与观测包括：
\begin{itemize}
    \item 宇宙学红移的时间漂移（time drift of cosmological redshifts）；\(^\text{[19]}\)
    \item 利用宇宙微波背景辐射（CMB）光子的反射来建模局部引力势（local gravitational potential）；\(^\text{[20]}\)
    \item 超新星光度的红移依赖关系（redshift dependence of supernova luminosity）；\(^\text{[21]}\)
    \item 与暗能量相关的运动性 Sunyaev–Zeldovich 效应（kinetic Sunyaev–Zeldovich effect）；\(^\text{[22]}\)
    \item 宇宙中微子背景（cosmic neutrino background）；\(^\text{[23]}\)
    \item 积分 Sachs–Wolfe 效应（integrated Sachs–Wolfe effect）；\(^\text{[24]}\)
    \item 对宇宙微波背景（CMB）各向同性与均匀性的检验；\(^\text{[25][26][27][28][29]}\)
    \item 对 KBC 空洞（KBC Void）的观测——部分研究者声称它违反了宇宙学原理与哥白尼原理，\(^\text{[30]}\)而另一些研究者则认为其与二者相容。\(^\text{[31]}\)
\end{itemize}
\subsection{无哥白尼原理的物理学}
标准宇宙学模型——即 $\Lambda$CDM 模型（Lambda-CDM model）——假设了哥白尼原理（Copernican principle）及其更一般的宇宙学原理（cosmological principle）。然而，一些宇宙学家与理论物理学家提出了不依赖这些原理的模型，旨在通过限制观测结果的取值范围、解决 $\Lambda$CDM 模型中已知的特定问题，并提出可用于区分现有模型与其他可能模型的实验检验方法。

在这一背景下，一个突出的例子是非均匀宇宙学（inhomogeneous cosmology），其目的是模拟观测到的宇宙加速膨胀以及宇宙常数（cosmological constant）现象。该类模型并不采用当前被广泛接受的“暗能量”（dark energy）概念，而是假设宇宙的非均匀性远比现有认知更强——例如，假设太阳系位于一个极其庞大的低密度空洞（low-density void）中。\(^\text{[32]}\)然而，为使模型符合观测结果，太阳系必须非常接近该空洞的中心，这一假设立即违背了哥白尼原理。

虽然宇宙学中的大爆炸模型（Big Bang model）有时被认为是由哥白尼原理与红移观测（redshift observations）共同推导而来，但即使在不假设哥白尼原理的情况下，大爆炸模型仍然可以被认为是有效的。因为宇宙微波背景（cosmic microwave background）、原初气体云（primordial gas clouds）、以及星系的结构、演化与分布等独立证据，都支持大爆炸理论的成立，而这些证据与哥白尼原理无关。然而，大爆炸模型中的关键命题——例如宇宙的膨胀——在此情形下本身也成为类似于哥白尼原理的假设，而非由哥白尼原理与观测共同导出。
\subsection{另见}
\begin{itemize}
\item 绝对空间与时间（Absolute space and time）
\item 人择原理（Anthropic principle）
\item “邪恶轴”现象（Axis of evil, cosmology） 
\item 哈勃气泡（Hubble Bubble, astronomy） 
\item 平庸原理（Mediocrity principle）
\item 粒子沙文主义（Particle chauvinism）
\item $P$ 对称性（P symmetry）
\item 稀有地球假说（Rare Earth hypothesis）
\item 纪录片《The Principle》（2014年）
\item 宇宙学原理（Cosmological principle）
\end{itemize}
\subsection{参考文献}
\begin{enumerate}
    \item Peacock, John A. (1998).Cosmological Physics. Cambridge University Press, p.~66. ISBN~978-0-521-42270-3.
    \item Bondi, Hermann (1952).Cosmology. Cambridge University Press, p.~13.
    \item Kuhn, Thomas S. (1957). The Copernican Revolution: Planetary Astronomy in the Development of Western Thought. Harvard University Press. Bibcode:1957crpa.book.....K. ISBN~978-0-674-17103-9.
    \item Musser, George (2001). “Copernican Counterrevolution.” Scientific American, 284(3):~24. Bibcode:2001SciAm.284c..24M. doi:10.1038/scientificamerican0301-24a.
    \item Danielson, Dennis (2009). “The Bones of Copernicus.” American Scientist, 97(1):~50–57. doi:10.1511/2009.76.50.
    \item Sagan, Carl. Cosmos. (1980), p.~193.
    \item Rowan-Robinson, Michael (1996). Cosmology (3rd ed.). Oxford University Press, pp.~62–63. ISBN~978-0-19-851884-6.
    \item Yadav, Jaswant; Bagla, J. S.; Khandai, Nishikanta (2010). “Fractal dimension as a measure of the scale of homogeneity.” Monthly Notices of the Royal Astronomical Society, 405(3):~2009–2015. arXiv:1001.0617. Bibcode:2010MNRAS.405.2009Y. doi:10.1111/j.1365-2966.2010.16612.x.
    \item Billings, Lee (2020). “Do We Live in a Lopsided Universe?” Scientific American. Retrieved March 24, 2022.
    \item Migkas, K. et al. (2020). “Probing cosmic isotropy with a new X-ray galaxy cluster sample through the LX–T scaling relation.” Astronomy \& Astrophysics, 636:~A15. arXiv:2004.03305. doi:10.1051/0004-6361/201936602.
    \item Secrest, Nathan J. et al.} (2021. “A Test of the Cosmological Principle with Quasars.” The Astrophysical Journal Letters, 908(2):~L51. arXiv:2009.14826. doi:10.3847/2041-8213/abdd40.
    \item Javanmardi, B. et al. (2015). “Probing the Isotropy of Cosmic Acceleration Traced by Type Ia Supernovae.”The Astrophysical Journal Letters, 810(1):~47. arXiv:1507.07560. doi:10.1088/0004-637X/810/1/47.
    \item Gott, J. Richard III et al. (2005). “A Map of the Universe.” The Astrophysical Journal, 624(2):~463–484. arXiv:astro-ph/0310571. doi:10.1086/428890.
    \item Horvath, I.; Hakkila, J.; Bagoly, Z. (2013). “The largest structure of the Universe, defined by Gamma-Ray Bursts.” arXiv:1311.1104 [astro-ph.CO].
    \item “Line of galaxies is so big it breaks our understanding of the universe.”
    \item Mazurenko, Sergi jet al.}(2024). “A Simultaneous Solution to the Hubble Tension and Observed Bulk Flow within 250 H$^{-1}$ Mpc.” Monthly Notices of the Royal Astronomical Society, 527(3):~4388–4396. doi:10.1093/mnras/stad3357; Haslbauer, Moritz; Banik, Indranil; Kroupa, Pavel (2020). “The KBC Void and Hubble Tension Contradict $\Lambda$CDM on a Gpc Scale – Milgromian Dynamics as a Possible Solution.” MNRAS, 499(2):~2845–2883. doi:10.1093/mnras/staa2348.
    \item Bondi, H.; Gold, T. (1948). “The Steady-State Theory of the Expanding Universe.” Monthly Notices of the Royal Astronomical Society, 108(3):~252–270. doi:10.1093/mnras/108.3.252.
    \item Clarkson, C.; Bassett, B.; Lu, T. (2008). “A General Test of the Copernican Principle.” Physical Review Letters, 101(1):~011301. arXiv:0712.3457. doi:10.1103/PhysRevLett.101.011301.
    \item Uzan, J. P.; Clarkson, C.; Ellis, G. (2008). “Time Drift of Cosmological Redshifts as a Test of the Copernican Principle.” Physical Review Letters, 100(19):~191303. arXiv:0801.0068. doi:10.1103/PhysRevLett.100.191303.
    \item Caldwell, R.; Stebbins, A. (2008). “A Test of the Copernican Principle.” Physical Review Letters, 100(19):~191302. arXiv:0711.3459. doi:10.1103/PhysRevLett.100.191302.
    \item Clifton, T.; Ferreira, P.; Land, K. (2008). “Living in a Void: Testing the Copernican Principle with Distant Supernovae.” Physical Review Letters, 101(13):~131302. arXiv:0807.1443. doi:10.1103/PhysRevLett.101.131302.
    \item Zhang, P.; Stebbins, A. (2011). “Confirmation of the Copernican Principle through the Anisotropic Kinetic Sunyaev–Zel’dovich Effect.” Philosophical Transactions of the Royal Society A, 369(1957):~5138–5145. doi:10.1098/rsta.2011.0294.
    \item Jia, J.; Zhang, H. (2008). “Can the Copernican Principle Be Tested Using the Cosmic Neutrino Background?”Journal of Cosmology and Astroparticle Physics, 2008(12):~002. arXiv:0809.2597. doi:10.1088/1475-7516/2008/12/002.
    \item Tomita, K.; Inoue, K. (2009). “Probing Violation of the Copernican Principle via the Integrated Sachs–Wolfe Effect.” Physical Review D, 79(10):~103505. arXiv:0903.1541. doi:10.1103/PhysRevD.79.103505.
    \item Clifton, T.; Clarkson, C.; Bull, P. (2012). “Isotropic Blackbody Cosmic Microwave Background Radiation as Evidence for a Homogeneous Universe.” Physical Review Letters, 109(5):~051303. arXiv:1111.3794. doi:10.1103/PhysRevLett.109.051303.
    \item Kim, J.; Naselsky, P. (2011). “Lack of Angular Correlation and Odd-Parity Preference in Cosmic Microwave Background Data.” The Astrophysical Journal, 739(2):~79. arXiv:1011.0377. doi:10.1088/0004-637X/739/2/79.
    \item Copi, C. J.; Huterer, D.; Schwarz, D. J.; Starkman, G. D. (2010). “Large-Angle Anomalies in the CMB.”Advances in Astronomy, 2010:~847541. arXiv:1004.5602. doi:10.1155/2010/847541.
    \item Ade, P. A. R. et al. (Planck Collaboration) (2013). “Planck 2013 results. XXIII. Isotropy and Statistics of the CMB.” Astronomy \& Astrophysics, 571:~A23. arXiv:1303.5083. doi:10.1051/0004-6361/201321534.
    \item Longo, Michael (2007). “Does the Universe Have a Handedness?” arXiv:astro-ph/0703325.
    \item Haslbauer, M.; Banik, I.; Kroupa, P. (2020). “The KBC Void and Hubble Tension Contradict $\Lambda$CDM on a Gpc Scale – Milgromian Dynamics as a Possible Solution.” MNRAS, 499(2):~2845–2883. doi:10.1093/mnras/staa2348.
    \item Sahlén, Martin; Zubeldía, Íñigo; Silk, Joseph (2016). “Cluster–Void Degeneracy Breaking: Dark Energy, Planck, and the Largest Cluster and Void.” The Astrophysical Journal Letters, 820(1):~L7. doi:10.3847/2041-8205/820/1/L7.
    \item February, S.; Larena, J.; Smith, M.; Clarkson, C. (2010). “Rendering Dark Energy Void.” Monthly Notices of the Royal Astronomical Society, 405(4):~2231. arXiv:0909.1479. doi:10.1111/j.1365-2966.2010.16627.x.
\end{enumerate}

